% Aula 05 --- as duas condições da Proposição da média, e o que a floresta faz com elas.
Justificamos o \emph{bagging} com uma proposição de enunciado curto: se
$g_1,\dots,g_B$ são não-viesados para $r(\x)$, têm a mesma variância $v(\x)$ e são
\textbf{não-correlacionados}, então a média $\bar g=\frac1B\sum_b g_b$ tem risco
$\Var(Y\mid\x)+v(\x)/B$, que nunca passa do risco $\Var(Y\mid\x)+v(\x)$ de cada
$g_b$.

Proposições com hipóteses só ficam úteis quando se sabe o que acontece ao violá-las
--- e, no caso do \emph{bagging}, uma delas é violada o tempo todo.
\begin{itens}
  \item \pts{1,1} Identifique as duas condições e explique, para cada uma, o que
        acontece com a conclusão se ela falhar.
  \item \pts{0,8} O \emph{bagging} usa árvores \textbf{profundas e não podadas}.
        Qual das duas condições essa escolha atende? E por que podar as árvores
        seria contraproducente aqui?
  \item \pts{0,9} Na prática as árvores do \emph{bagging} são correlacionadas.
        Explique a razão --- o que elas têm em comum que as faz parecidas --- e
        por que aumentar $B$ não resolve esse problema específico.
  \item \pts{1,2} A floresta aleatória sorteia $m<p$ covariáveis candidatas em
        cada nó. Explique como isso ataca a condição comprometida, qual é o preço
        pago na outra, e por que essa troca costuma compensar.
\end{itens}

\begin{solucao}
\textbf{(a)} As duas condições são:

\textbf{(i) não-viesados.} Se os $g_b$ tiverem viés, a média o herda
integralmente: $\E[\bar g]=\E[g_b]\ne r(\x)$. Promediar reduz variância, mas
\textbf{não corrige viés} --- a média de mil estimadores enviesados é igualmente
enviesada. A conclusão sobre o risco deixa de valer, porque sobra um termo de
viés$^2$ que $B$ nenhum elimina.

\textbf{(ii) não-correlacionados.} Se houver correlação $\rho>0$ entre pares, a
variância da média não é $v/B$, e sim
\[
  \Var(\bar g)=\rho\,v+\frac{1-\rho}{B}\,v \;\xrightarrow[B\to\infty]{}\; \rho\,v .
\]
A redução \textbf{estaciona num piso} $\rho v$, em vez de ir a zero.

\smallskip
\textbf{(b)} Atende a condição \textbf{(i)}, a de não-viesamento.

Árvores profundas são muito flexíveis: acompanham $r$ de perto e têm viés baixo,
ao custo de variância alta. E variância alta é justamente o defeito que a
agregação sabe consertar --- de graça, pela própria proposição.

Podar seria contraproducente porque a poda \emph{introduz viés de propósito} para
reduzir variância. Só que a variância já vai ser tratada pela média dos $B$
ajustes. Podar seria, então, pagar viés por um benefício que se teria sem pagar
nada.

A regra que fica: deixe cada árvore ser instável, e deixe a média estabilizar.

\smallskip
\textbf{(c)} As árvores do \emph{bagging} são treinadas em amostras bootstrap do
\textbf{mesmo} conjunto de dados, e todas enxergam \textbf{todas} as $p$
covariáveis em todos os nós.

Se houver uma covariável fortemente preditiva, ela tende a ser escolhida no
primeiro corte por quase todas as árvores. E árvores que começam igual continuam
parecidas: elas erram juntas, o que é precisamente o que $\rho>0$ significa.

Aumentar $B$ não resolve porque, na fórmula de (a), o termo $\rho v$ \textbf{não
depende de $B$}. Acrescentar árvores só encolhe o segundo termo; o piso permanece
onde está. Passado certo ponto, cada árvore nova custa tempo de computação e não
compra quase nada.

\smallskip
\textbf{(d)} A floresta ataca diretamente o $\rho$. Ao sortear $m<p$ candidatas
por nó, ela impede que a covariável dominante seja usada sempre: em cerca de
$1-m/p$ dos nós ela sequer está disponível, e a árvore é obrigada a olhar para
outra coisa. As árvores passam a divergir estruturalmente, $\rho$ cai, e com ele o
piso $\rho v$.

O preço é pago na qualidade individual. Uma árvore proibida de usar a melhor
covariável em parte dos nós é uma árvore \textbf{pior}, com variância individual
$v$ maior. O sorteio degrada cada membro para melhorar o conjunto --- o que é
contraintuitivo, e é o ponto que costuma ser mal explicado.

Compensa porque, com $B$ grande, o que sobra em $\rho v+\frac{1-\rho}{B}v$ é
$\rho v$, e a queda em $\rho$ costuma ser proporcionalmente maior que a subida em
$v$. Há um limite, claro: com $m$ pequeno demais as árvores ficam ruins demais e o
produto $\rho v$ volta a subir. É por isso que $m$ é um hiperparâmetro, e
$m\approx p/3$ em regressão é um ponto de partida razoável --- não uma verdade.
\end{solucao}
